\documentclass[prd,twocolumn,aps,nobibnotes,superscriptaddress,preprintnumbers,floatfix]{revtex4-1}
\bibliographystyle{apsrev4-1}
\usepackage{graphicx,bm,amsmath,amssymb,xcolor}
\usepackage{hyperref}
\graphicspath{{./images/}}

%%% R9: the letter recommended by both round-2 referees. Its scientific content is
%%% Sec.~\ref{sec:magnetic-uetc} of derivation.tex (the companion), which supplies the
%%% derivation chain, the kernel definitions and every cross-check quoted here.

\begin{document}

\title{No usable cusp in the magnetic stress correlator:\\
       the turbulent gravitational-wave peak is a window effect}

\author{Murman Gurgenidze}
\email{mgurgeni@andrew.cmu.edu}
\affiliation{McWilliams Center for Cosmology and Department of Physics,
Carnegie Mellon University, Pittsburgh, PA 15213, USA}

\author{Tina Kahniashvili}
\email{tinatin@andrew.cmu.edu}
\affiliation{McWilliams Center for Cosmology and Department of Physics,
Carnegie Mellon University, Pittsburgh, PA 15213, USA}
\affiliation{School of Natural Sciences and Medicine, Ilia State University, 0194 Tbilisi, Georgia}
\affiliation{Abastumani Astrophysical Observatory, Tbilisi, GE-0179, Georgia}

\begin{abstract}
The gravitational-wave spectrum radiated by primordial turbulence is fixed, at fixed
spatial source, by the unequal-time correlator (UETC) of the anisotropic stress. Whether
that correlator is differentiable at zero time lag decides the shape of the spectrum
above its peak: a cusp gives a heavy $\omega^{-2}$ tail and displaces the peak to the
source scale, while a smooth correlator does neither. The velocity UETC of hydrodynamic
turbulence has been measured and is Gaussian, hence cuspless; the corresponding magnetic
\emph{stress} correlator, which is what actually sources the waves, has not been. We show
that it can be constrained without any two-time data. In a magnetohydrodynamic simulation
started from rest, the gravitational-wave field is itself a running cosine transform of
the stress UETC evaluated on the light cone, so the radiated energy at the end of the run
is a parameter-free \emph{prediction} for each candidate correlator once the stress and
the Alfv\'en speed are read off. Applied to a published run, cuspless correlators
reproduce the radiated energy while cusped ones overpredict it by one to two orders of
magnitude. The magnetic stress correlator is therefore no more cusped than the
hydrodynamic one, the $\omega^{-2}$ tail is excluded for this class of source, and the
source-scale peak seen in simulations is produced by the finite emission window rather
than by decorrelation. The method needs only quantities such simulations already write to
disk.
\end{abstract}

\maketitle

\section{Introduction}
%%% ~1 pp. From derivation.tex Sec. I, compressed: the three inputs; that the temporal
%%% input is the least constrained; that the cusp/no-cusp question is binary and decides
%%% the UV and the peak. Cite Kraichnan:1964, Gogoberidze:2007an, Caprini:2006jb,
%%% Caprini:2009fx, Niksa:2018ofa, Auclair:2022jod, RoperPol:2019wvy, RoperPol:2022iel.
%%% State plainly what is already known and what is new: the Abelian tail result and its
%%% peak consequence are Caprini:2009fx's; what is new is the measurement.

\section{The criterion}
%%% ~0.75 pp. eq:cusp-tail, the two-leg rule (quadratic source carries the SQUARED
%%% correlator), and the collapse: every cusped model gives the same omega^-2, so the
%%% question is binary. Note the window/lag-memory distinction and the coefficients
%%% 2.33/W vs pi/tau_c -- see companion Sec. V.B.

\section{The measurement}
\label{sec:letter-measurement}
%%% ~1.5 pp -- the core. From derivation.tex Sec. sec:magnetic-uetc verbatim:
%%%  (a) why a direct UETC is impossible from published data: cadence 1e-3 vs the lag
%%%      1/k ~ 8e-4 at the GW peak;
%%%  (b) the identity that rescues it: with h = hdot = 0, ghat(k,t) = int cos(k(t-t')) S;
%%%  (c) the parameter-free forward model and its table;
%%%  (d) robustness: zero-padding, v_A at start/mean/end;
%%%  (e) the three limits -- Gaussian closure, asymmetric discrimination (coherent and
%%%      Gaussian sweeping are not separated), one Pm=1 run with an imposed initial field.
\begin{table}[b]
  \centering
  \begin{tabular}{@{}lc@{}}
  \hline\hline
  correlator & $E_{\rm model}/E_{\rm data}$ \\
  \hline
  coherent, $f=1$ & $0.35$--$1.19$ \\
  Gaussian sweeping, $f=e^{-(kv_A\tau)^2/2}$ & $0.90$--$1.14$ \\
  cusped, $f=e^{-kv_A|\tau|}$ & $19.5$--$97.8$ \\
  BK2016, $f=(1+kv_A\tau)^{-2/3}$ & $13.0$--$64.0$ \\
  \hline\hline
  \end{tabular}
  \caption{Parameter-free forward models of the radiated energy over $p=1.5$--$30$.}
  \label{tab:letter-uetc}
\end{table}

\section{Consequences}
%%% ~1 pp.
%%%  - for templates: drop power-law decorrelation from turbulence templates;
%%%  - for the peak: it is a window effect, so it is RoperPol:2022iel's constant-source
%%%    mechanism, and we say so;
%%%  - what would overturn this: a second run, a driven rather than imposed field, Pm != 1.

\section{Summary}
%%% ~0.25 pp. Point at the companion for the kernel compendium, the branch diagram, the
%%% break in Hz, and the dissipation outlook.

\bibliography{bib}
\end{document}
